\section{Use for the comparison between group and solo flight}\label{sec:comparison}

\paragraph{Labels on the phases.} For a run $\sigma$ of one phase (climb, search or glide)
of flight $i$, with interval $J_\sigma$, define the fraction in a group and the mean size:
\begin{equation}\label{eq:segment}
  f_\sigma=\frac{\bigl|\{t\in J_\sigma\cap T^\dep_i:\ \gamma_i(t)\notin\{\st{L},\st{S}\}\}\bigr|}{|J_\sigma\cap T^\dep_i|},
  \qquad
  \bar n_\sigma=\text{mean of } n_i(t) \text{ over } J_\sigma\cap T^\dep_i ,
\end{equation}
with $n_i(t)$ the number of members of the configuration of $i$ at $t$, 1 if solo. The run is
\emph{in a group} if $f_\sigma\ge1/2$ and \emph{solo} if $f_\sigma=0$; the intermediate cases are
excluded, or $f_\sigma$ is used as a continuous covariate (Figure~\ref{fig:segments}). Runs in the
local phase are excluded.

\begin{figure}
\centering
\begin{tikzpicture}[x=1.1cm,y=1cm,
    seg/.style={draw=white, line width=1pt, minimum height=6mm, inner sep=0pt, font=\footnotesize}]
  \node[lbl, anchor=east] at (-0.1,1.2) {$\gamma_i(t)$};
  \fill[light] (0,1.05) rectangle (3.0,1.35);
  \fill[blueI] (3.0,1.05) rectangle (9.3,1.35);
  \fill[light] (9.3,1.05) rectangle (12,1.35);
  \node[font=\footnotesize] at (1.5,1.2) {\st{S}};
  \node[font=\footnotesize, text=white] at (6.15,1.2) {configuration};
  \node[font=\footnotesize] at (10.65,1.2) {\st{S}};
  \node[lbl, anchor=east] at (-0.1,0.3) {phases};
  \foreach \a/\b/\name/\f/\cls in {0/2/climb/0/solo, 2/4.5/glide/0.6/group, 4.5/6.5/climb/1/group,
                                   6.5/9/glide/1/group, 9/10/search/0.3/excluded, 10/12/climb/0/solo} {
    \fill[ink2!25] (\a,0.0) rectangle (\b,0.6);
    \draw[white, line width=1.5pt] (\a,0) -- (\a,0.6);
    \node[font=\footnotesize] at ({(\a+\b)/2},0.3) {\name};
    \node[font=\footnotesize, anchor=north] at ({(\a+\b)/2},-0.05) {$f_\sigma=\f$};
    \node[font=\footnotesize, anchor=north, text=ink2] at ({(\a+\b)/2},-0.45) {\cls};
  }
  \draw[->, ink2] (0,-1.05) -- (12.3,-1.05) node[right, lbl, ink] {$t$};
\end{tikzpicture}
\caption{Labels on the phases of one flight. Each run of one phase gets the fraction
$f_\sigma$ of its duration spent in a configuration: runs with $f_\sigma\ge1/2$ count as group,
those with $f_\sigma=0$ as solo, the search with $f_\sigma=0.3$ is left out.}
\label{fig:segments}
\end{figure}

\paragraph{Stratification.} Group flights concentrate on the good days and the busy sites: a
global comparison would mix the effect of the group with that of the weather. The comparison runs
within strata $s$ (date and area):
\begin{equation}\label{eq:strata}
  \hat\Delta=\sum_s w_s\bigl(\hat\theta_{G,s}-\hat\theta_{S,s}\bigr),
\end{equation}
where $\hat\theta$ is the statistic of interest and the weights $w_s$ are proportional, for
example, to the number of group runs of the stratum.

\paragraph{Uncertainties.} Runs of pilots of the same episode are not independent, and neither
are those of the same flight. The bootstrap resamples the dependency blocks
(Section~\ref{sec:blocks}) within the strata.

\paragraph{Confounders.} Competition days; the level of the pilot (the same pilots often fly
together); the glider class (EN A--D); the time of day.

\paragraph{Questions to ask.} Distributions of the times spent in thermals and of the lengths of
the glides, in a group and alone; mean squared displacement and effective exponent; differences
between followers and leaders, by the sign of $\ell_{\rm med}$. In birds, followers leave the
thermals earlier and lower~\cite{flack2018}, and individuals with social information glide
faster~\cite{williams2018}: these are hypotheses to test on pilots before assuming them.
